\section{Анализ существующих исследований}

\subsection{Распределение модели и ограничения памяти}

Система \emph{Petals} \cite{petals2023} позволяет исполнять последовательные блоки больших языковых моделей на узлах нескольких участников. Она демонстрирует практически значимый сценарий распределения слоёв по сети, но включает задачи динамической маршрутизации, управления участниками и отказоустойчивости. В предлагаемой работе изучается более контролируемый стенд: исполнители фиксированы, а основной предмет анализа~--- стоимость обмена между стадиями и воспроизводимость измерений.

\emph{DeepSpeed Inference} \cite{deepspeed2022} объединяет методы параллелизма и размещения вычислений для крупных Transformer-моделей в многоускорительной среде. Исследование показывает, что оптимальная организация инференса зависит от ограничивающего ресурса и выбранной нагрузки. Непосредственное сравнение скорости DeepSpeed с локальным CPU-прототипом не имеет смысла из-за разницы масштабов, однако принципы моделирования ограничений являются релевантными.

\emph{FlexGen} \cite{flexgen2023} ориентирован на генерацию при ограниченной GPU-памяти. В нём совместно используются GPU, RAM и накопитель, а выбор стратегии размещения данных формулируется как задача оптимизации. В отличие от FlexGen, здесь предполагается оптимизировать границы \emph{последовательных групп Transformer-слоёв} и анализировать стоимость межпроцессного обмена в Ray. Идея учёта ресурсоёмкости отдельных компонентов при этом концептуально сходна.

\subsection{Обслуживание запросов, коммуникации и балансировка}

В системе \emph{vLLM} \cite{vllm2023} механизм PagedAttention применяется для управления KV-кэшем при множестве запросов. Система \emph{DistServe} \cite{distserve2024} исследует раздельное размещение фаз обработки входного контекста и декодирования. Обе работы показывают значимость формы нагрузки и состояния вычислительного процесса; текущая версия нашего прототипа не содержит KV-кэш и многозапросную конвейеризацию.

Jekel и Bhatele \cite{jekel2026} исследуют коммуникационные ограничения многоузлового GPU-инференса, уделяя особое внимание коллективным операциям при тензорном параллелизме. Их задача отличается от последовательного обмена активациями между Ray Actors, поэтому численные результаты нельзя переносить на локальную CPU-систему. Тем не менее работа обосновывает необходимость раздельного анализа вычислительных и коммуникационных компонентов задержки.

Hwang и Ahn \cite{hwang2026} исследуют балансировку стадий конвейерного обслуживания LLM и показывают, что дисбаланс и простои способны ограничивать пропускную способность. Это подтверждает значимость оптимизации размещения слоёв, но подчёркивает необходимость различать два режима: \emph{последовательную обработку одного запроса} и \emph{перекрывающийся конвейер нескольких запросов}.

Фреймворк \emph{Ray} \cite{ray2018} обеспечивает удалённые объекты и планирование по логическим ресурсам. Выбор Ray является инженерным решением для прототипа, а не предметом заявляемой научной новизны. Логические ресурсы Ray не обеспечивают жёсткой привязки к отдельным физическим ядрам \cite{rayresources}.

\begin{table*}[t]
\centering\small
\caption{Сопоставление работ и исследовательской постановки}
\label{tab:related}
\begin{tabularx}{\textwidth}{@{}p{2.55cm}X X@{}}
\toprule
\textbf{Подход} & \textbf{Основной предмет} & \textbf{Связь с данной работой}\\
\midrule
Petals \cite{petals2023} & Размещение блоков LLM по сети & Послойный обмен; здесь фиксированные исполнители\\
DeepSpeed \cite{deepspeed2022} & Инференс крупных моделей на GPU & Учёт вычислительных и ресурсных ограничений\\
FlexGen \cite{flexgen2023} & Размещение данных при дефиците памяти & Оптимизационная постановка при ограниченных ресурсах\\
vLLM, DistServe \cite{vllm2023,distserve2024} & KV-кэш и организация обслуживания & Ориентиры для будущего расширения; нет в исходной системе\\
Jekel--Bhatele \cite{jekel2026} & Коммуникации многоузлового GPU-инференса & Обоснование декомпозиции накладных расходов\\
Hwang--Ahn \cite{hwang2026} & Балансировка конвейерных стадий & Различие между задержкой и конвейерной пропускной способностью\\
Настоящая работа & Сериализация выходных данных и выбор разбиения & Воспроизводимая CPU-постановка с проверяемыми ограничениями\\
\bottomrule
\end{tabularx}
\end{table*}

Обзор выявляет конкретную нишу исследования: на простом управляемом стенде последовательно выделить эффект объёма возвращаемых логитов, стоимость их вычисления и влияние выбора границ разбиения на прогнозируемую задержку. Такой подход не заменяет промышленного сервинга, но позволяет получить воспроизводимые результаты о свойствах исследуемой архитектуры.
